\section{Conclusão}
\label{sec:conclusao}

Este relatório consolidou o Trabalho 1 da disciplina de Redes
Neurais Sem Peso (mestrado, 2026/1), em que reproduzimos
fielmente os oito \emph{baselines} de Alsaedi et al.\
(2020)~\cite{alsaedi2020toniot} no dataset TON\_IoT e avaliamos cinco
modelos da família WiSARD sob \emph{exatamente o mesmo protocolo}.

\paragraph{Quanto à reprodução do paper.} Sob o protocolo
\emph{paper-faithful} (\emph{split} 80/20, $k=4$ \emph{StratifiedKFold},
métrica \texttt{F1 weighted}), nossa reprodução dos oito
\emph{baselines} bate com as Tabelas 10--11 do paper dentro de
$|\Delta|<0{,}15$ em 91\% dos pares (modelo, base), com média de
$\Delta = -0{,}009$. Os deltas residuais concentram-se nas bases mais
afetadas pela divergência de tamanho entre os CSVs públicos e o
dataset original do paper (sub-amostragem da classe normal para 43\%
do volume original), o que documentamos como o fator limitante
dominante da reprodução.

\paragraph{Quanto à família WiSARD.} Sob o mesmo protocolo, a família
WiSARD apresenta desempenho competitivo na maioria das bases
\emph{per-device}: WiSARD ($0{,}824$) e ClusWiSARD ($0{,}822$) lideram o
grupo \emph{praticamente empatadas}, a cerca de um ponto percentual de
Random Forest e CART ($\approx 0{,}83$) e à frente de kNN ($0{,}789$),
LSTM ($0{,}712$) e SVM ($0{,}659$); BloomWiSARD ($0{,}764$), BTHOWeN
($0{,}760$) e ULEEN ($0{,}758$) vêm logo atrás. Essa paridade
WiSARD/ClusWiSARD só emergiu após unificarmos os tetos de capacidade do
\emph{sweep} (\S{}\ref{sec:wisard_fair}): com tetos desiguais, ClusWiSARD
aparecia $0{,}04$ atrás por artefato de cobertura, não de modelo. Em
bases com \emph{leak} categórico forte (\texttt{Fridge},
\texttt{Garage\_Door}), toda a família satura em F1 $\approx 1{,}000$;
em bases sem sinal discriminativo (\texttt{Motion\_Light}), todos
colapsam para predição majoritária. Em \texttt{Modbus}, com estrutura
não-linear nos códigos FC, WiSARD e ClusWiSARD com endereçamento amplo
chegam a $0{,}95$, \emph{igualando} os modelos de árvore (RF/CART
$\approx 0{,}94$). No \emph{combined} multi-classe, o \emph{gap} cresce:
a família fica em $0{,}24$--$0{,}43$, com ULEEN ($0{,}425$), WiSARD e
ClusWiSARD ($0{,}413$) à frente de LR/LDA/SVM ($\approx 0{,}30$) e bem
acima de NB ($0{,}03$), mas distante de RF/CART ($0{,}60$).

\paragraph{Quanto à viabilidade em \emph{edge}.} A discussão mais
relevante deste trabalho não está nos números brutos de F1, mas no
\emph{trade-off} F1 $\times$ custo, e nesse \emph{trade-off} o eixo
decisivo é a memória. O BTHOWeN treina em $0{,}09$\,s medianos (cerca de
\textbf{705$\times$ mais rápido que o LSTM}) e entrega F1 $\approx 0{,}76$.
Em F1 $\times$ tempo de CPU, porém, os modelos de árvore ainda dominam:
RF e kNN treinam tão rápido quanto o BTHOWeN e com F1 superior. A
vantagem genuína da família aparece na memória, sob uma métrica única de
estado implantado: as variantes de filtro (BloomWiSARD, BTHOWeN, ULEEN)
são implantáveis em $0{,}75$--$2$\,\emph{KB} \emph{fixos}, cerca de três
ordens de magnitude abaixo dos modelos de árvore (RF $\approx 3$\,MB,
kNN $\approx 1$\,MB), a um custo de poucos pontos de F1; WiSARD e
ClusWiSARD recuperam o melhor F1 da família ($0{,}82$) com
$\approx 2$\,KB \emph{medianos}, ainda muito abaixo das árvores, mas
com \emph{footprint} dependente dos dados que cresce nas bases
estruturadas. É esse \emph{footprint} reduzido (somado à inferência em
\emph{hardware} dedicado documentada nos artigos originais de BTHOWeN e
ULEEN, que um \emph{benchmark} em CPU não captura) que torna a
família preferível para \emph{deployment} em hardware embarcado com
restrição de memória.

\paragraph{Trabalhos futuros.} As direções naturais para este trabalho
são (i) acesso ao dataset original do paper, eliminando a divergência
de tamanho que limita a reprodução; (ii) \emph{dispatch} por \emph{feature}, escolhendo o melhor
termômetro para cada variável em vez de aplicar um único termômetro
uniformemente, uma vez que a alocação \texttt{non-uniform} já entrou no
\emph{sweep} atual;
(iii) avaliação do BTHOWeN em \emph{ULEEN-style training}
(\emph{epochs}+\emph{learning rate}) para verificar se o tempo extra
de treinamento compensa pelo F1 ganho; (iv) extensão dos
\emph{sweeps} para hardware-target (latência de inferência e
\emph{memory footprint} sob restrição de FPGA), seguindo o protocolo
do BTHOWeN original.

\paragraph{Conclusão final.} A família WiSARD é uma família viável de
modelos sem peso para detecção de anomalia em IoT, com perfil
F1 $\times$ custo claramente distinto e em geral favorável frente aos
\emph{baselines} clássicos para o cenário \emph{edge}. Os resultados
não substituem os modelos baseados em árvore quando F1 puro é o
critério único, mas oferecem redução de cerca de três ordens de
magnitude em \emph{footprint} de memória (variantes de filtro, com
\emph{footprint} fixo) a uma perda relativamente pequena em F1, nas
bases sem leak forte. Duas contribuições metodológicas acompanham esse
resultado: a unificação dos tetos de capacidade do \emph{sweep}, que
tornou a comparação intra-família justa (revelando que a aparente
liderança da WiSARD sobre a ClusWiSARD era artefato de cobertura de
\emph{grid}), e a métrica única de estado implantado que coloca todos
os modelos sem peso no mesmo eixo de memória. A maior contribuição,
porém, é o diagnóstico dos problemas estruturais dos CSVs públicos do
TON\_IoT: o \emph{leak} categórico e a divergência de tamanho contra o
paper original, não discutidos no paper, e a caracterização do
\emph{leak} temporal (que o paper já descarta por risco de
\emph{overfit}) como determinístico. Esperamos que sirva de referência
para trabalhos futuros sobre esse dataset.

\paragraph{Disponibilidade de código e dados.} Todo o pipeline (o
\emph{notebook} de referência \texttt{06\_consolidado\_grupos.ipynb},
os \emph{scripts} de \emph{sweep} e \emph{benchmark} e os JSONL/JSON de
resultados que sustentam cada número reportado) está disponível em
\url{https://github.com/muanlartins/masters/tree/toniot-wisard-repro}.
A biblioteca \texttt{wisardpkg} (com os cinco termômetros ajustados
\emph{Distributive}, \emph{Gaussian}, \emph{Exponential},
\emph{Logarithmic} e \emph{Non-uniform}, e a métrica
\texttt{deployedSizeBytes}) é fixada por \emph{commit} em
\texttt{requirements.txt}. O dataset TON\_IoT
\texttt{Train\_Test\_IoT\_*.csv} é público (UNSW Canberra Cyber).

\paragraph{Uso de assistentes de IA.} O texto deste relatório foi
redigido com assistência de IA generativa (Claude Code), a partir
dos números, figuras e estrutura definidos pelos autores no
\emph{notebook} de consolidação. Toda a metodologia, todos os
experimentos e todas as conclusões são de autoria humana.
